\hspace{3mm}

\centerline{\bf \Huge Страшное слово на букву "И"}

\begin{flushright}
    \scriptsize{}
\end{flushright}

\problem{1} 
Даны натуральные числа $x_1, ..., x_n$. Докажите, что число 
$(1+x_1^2)(1+x_2^2)\ldots(1+x_n^2)$
можно представить в виде суммы квадратов двух целых чисел.

\problem{2} 
Проведём в выпуклом многоугольнике некоторые диагонали так, что никакие две из них не пересекаются (из одной вершины могут выходить несколько диагоналей). Доказать, что найдутся по крайней мере две вершины многоугольника, из которых не проведено ни одной диагонали.

\problem{3}
В прямоугольнике $3\times n$ стоят фишки трёх цветов, по n штук каждого цвета. Доказать, что можно переставить фишки в каждой строке так, чтобы в каждом столбце были фишки всех цветов.

\problem{4} 
Несколько прямых делят плоскость на части. Докажите, что эти части можно раскрасить в 2 цвета так, что граничащие части будут иметь разный цвет.

\problem{5} 
a) На сколько частей делят плоскость n прямых общего положения, то есть таких, что никакие две не параллельны и никакие три не проходят через одну точку?

\noindentб) На какое наибольшее число частей могут разбить плоскость n окружностей?




